\documentclass[11pt]{article}
% Préambule commun aux cours de la formation C++.
% Inclus par chaque cours.tex via % Préambule commun aux cours de la formation C++.
% Inclus par chaque cours.tex via % Préambule commun aux cours de la formation C++.
% Inclus par chaque cours.tex via \input{preambule}.

\usepackage[utf8]{inputenc}
\usepackage[T1]{fontenc}
\usepackage[french]{babel}
\usepackage{lmodern}
\usepackage{microtype}

\usepackage{geometry}
\geometry{a4paper, margin=2.5cm}

\usepackage{xcolor}
\usepackage{amsmath, amssymb}
\usepackage{enumitem}
\usepackage{hyperref}
\hypersetup{
    colorlinks=true,
    linkcolor=blue!60!black,
    urlcolor=blue!60!black,
    citecolor=blue!60!black,
}

% --- Coloration syntaxique du C++ via listings -----------------------------
\usepackage{listings}

\definecolor{codebg}{rgb}{0.97,0.97,0.95}
\definecolor{codekw}{rgb}{0.10,0.10,0.60}
\definecolor{codecomment}{rgb}{0.20,0.50,0.20}
\definecolor{codestring}{rgb}{0.65,0.15,0.15}
\definecolor{coderule}{rgb}{0.80,0.80,0.80}

% Permet les caractères accentués (UTF-8) à l'intérieur des blocs de code.
\lstset{
    literate=%
        {é}{{\'e}}1 {è}{{\`e}}1 {ê}{{\^e}}1 {ë}{{\"e}}1
        {à}{{\`a}}1 {â}{{\^a}}1 {ä}{{\"a}}1
        {î}{{\^i}}1 {ï}{{\"i}}1
        {ô}{{\^o}}1 {ö}{{\"o}}1
        {ù}{{\`u}}1 {û}{{\^u}}1 {ü}{{\"u}}1
        {ç}{{\c c}}1 {É}{{\'E}}1 {È}{{\`E}}1 {À}{{\`A}}1
}

\lstdefinestyle{cpp}{
    language=C++,
    backgroundcolor=\color{codebg},
    basicstyle=\ttfamily\small,
    keywordstyle=\color{codekw}\bfseries,
    commentstyle=\color{codecomment}\itshape,
    stringstyle=\color{codestring},
    showstringspaces=false,
    numbers=left,
    numberstyle=\tiny\color{gray},
    numbersep=8pt,
    frame=single,
    rulecolor=\color{coderule},
    breaklines=true,
    tabsize=4,
    morekeywords={constexpr,noexcept,override,final,nullptr,auto,
        std,string,vector,size_t,uint64_t,int64_t,optional,namespace}
}
\lstset{style=cpp}

% Raccourci pour du code en ligne, en police machine à écrire.
\newcommand{\code}[1]{\texttt{#1}}

% Boîte « À retenir »
\usepackage{tcolorbox}
\newtcolorbox{aretenir}[1][]{
    colback=yellow!8, colframe=orange!70!black,
    title=#1, fonttitle=\bfseries, sharp corners, boxrule=0.5pt
}
\newtcolorbox{piege}[1][Piège]{
    colback=red!5, colframe=red!60!black,
    title=#1, fonttitle=\bfseries, sharp corners, boxrule=0.5pt
}
.

\usepackage[utf8]{inputenc}
\usepackage[T1]{fontenc}
\usepackage[french]{babel}
\usepackage{lmodern}
\usepackage{microtype}

\usepackage{geometry}
\geometry{a4paper, margin=2.5cm}

\usepackage{xcolor}
\usepackage{amsmath, amssymb}
\usepackage{enumitem}
\usepackage{hyperref}
\hypersetup{
    colorlinks=true,
    linkcolor=blue!60!black,
    urlcolor=blue!60!black,
    citecolor=blue!60!black,
}

% --- Coloration syntaxique du C++ via listings -----------------------------
\usepackage{listings}

\definecolor{codebg}{rgb}{0.97,0.97,0.95}
\definecolor{codekw}{rgb}{0.10,0.10,0.60}
\definecolor{codecomment}{rgb}{0.20,0.50,0.20}
\definecolor{codestring}{rgb}{0.65,0.15,0.15}
\definecolor{coderule}{rgb}{0.80,0.80,0.80}

% Permet les caractères accentués (UTF-8) à l'intérieur des blocs de code.
\lstset{
    literate=%
        {é}{{\'e}}1 {è}{{\`e}}1 {ê}{{\^e}}1 {ë}{{\"e}}1
        {à}{{\`a}}1 {â}{{\^a}}1 {ä}{{\"a}}1
        {î}{{\^i}}1 {ï}{{\"i}}1
        {ô}{{\^o}}1 {ö}{{\"o}}1
        {ù}{{\`u}}1 {û}{{\^u}}1 {ü}{{\"u}}1
        {ç}{{\c c}}1 {É}{{\'E}}1 {È}{{\`E}}1 {À}{{\`A}}1
}

\lstdefinestyle{cpp}{
    language=C++,
    backgroundcolor=\color{codebg},
    basicstyle=\ttfamily\small,
    keywordstyle=\color{codekw}\bfseries,
    commentstyle=\color{codecomment}\itshape,
    stringstyle=\color{codestring},
    showstringspaces=false,
    numbers=left,
    numberstyle=\tiny\color{gray},
    numbersep=8pt,
    frame=single,
    rulecolor=\color{coderule},
    breaklines=true,
    tabsize=4,
    morekeywords={constexpr,noexcept,override,final,nullptr,auto,
        std,string,vector,size_t,uint64_t,int64_t,optional,namespace}
}
\lstset{style=cpp}

% Raccourci pour du code en ligne, en police machine à écrire.
\newcommand{\code}[1]{\texttt{#1}}

% Boîte « À retenir »
\usepackage{tcolorbox}
\newtcolorbox{aretenir}[1][]{
    colback=yellow!8, colframe=orange!70!black,
    title=#1, fonttitle=\bfseries, sharp corners, boxrule=0.5pt
}
\newtcolorbox{piege}[1][Piège]{
    colback=red!5, colframe=red!60!black,
    title=#1, fonttitle=\bfseries, sharp corners, boxrule=0.5pt
}
.

\usepackage[utf8]{inputenc}
\usepackage[T1]{fontenc}
\usepackage[french]{babel}
\usepackage{lmodern}
\usepackage{microtype}

\usepackage{geometry}
\geometry{a4paper, margin=2.5cm}

\usepackage{xcolor}
\usepackage{amsmath, amssymb}
\usepackage{enumitem}
\usepackage{hyperref}
\hypersetup{
    colorlinks=true,
    linkcolor=blue!60!black,
    urlcolor=blue!60!black,
    citecolor=blue!60!black,
}

% --- Coloration syntaxique du C++ via listings -----------------------------
\usepackage{listings}

\definecolor{codebg}{rgb}{0.97,0.97,0.95}
\definecolor{codekw}{rgb}{0.10,0.10,0.60}
\definecolor{codecomment}{rgb}{0.20,0.50,0.20}
\definecolor{codestring}{rgb}{0.65,0.15,0.15}
\definecolor{coderule}{rgb}{0.80,0.80,0.80}

% Permet les caractères accentués (UTF-8) à l'intérieur des blocs de code.
\lstset{
    literate=%
        {é}{{\'e}}1 {è}{{\`e}}1 {ê}{{\^e}}1 {ë}{{\"e}}1
        {à}{{\`a}}1 {â}{{\^a}}1 {ä}{{\"a}}1
        {î}{{\^i}}1 {ï}{{\"i}}1
        {ô}{{\^o}}1 {ö}{{\"o}}1
        {ù}{{\`u}}1 {û}{{\^u}}1 {ü}{{\"u}}1
        {ç}{{\c c}}1 {É}{{\'E}}1 {È}{{\`E}}1 {À}{{\`A}}1
}

\lstdefinestyle{cpp}{
    language=C++,
    backgroundcolor=\color{codebg},
    basicstyle=\ttfamily\small,
    keywordstyle=\color{codekw}\bfseries,
    commentstyle=\color{codecomment}\itshape,
    stringstyle=\color{codestring},
    showstringspaces=false,
    numbers=left,
    numberstyle=\tiny\color{gray},
    numbersep=8pt,
    frame=single,
    rulecolor=\color{coderule},
    breaklines=true,
    tabsize=4,
    morekeywords={constexpr,noexcept,override,final,nullptr,auto,
        std,string,vector,size_t,uint64_t,int64_t,optional,namespace}
}
\lstset{style=cpp}

% Raccourci pour du code en ligne, en police machine à écrire.
\newcommand{\code}[1]{\texttt{#1}}

% Boîte « À retenir »
\usepackage{tcolorbox}
\newtcolorbox{aretenir}[1][]{
    colback=yellow!8, colframe=orange!70!black,
    title=#1, fonttitle=\bfseries, sharp corners, boxrule=0.5pt
}
\newtcolorbox{piege}[1][Piège]{
    colback=red!5, colframe=red!60!black,
    title=#1, fonttitle=\bfseries, sharp corners, boxrule=0.5pt
}


\title{\textbf{Chapitre 3 — Lire le monde : I/O, flux et parsing}\\[0.3em]
       \large Entrées/sorties en profondeur et conversion de texte}
\author{Formation C++ moderne}
\date{}

\begin{document}
\maketitle

\begin{abstract}
\noindent Un programme utile parle au monde extérieur : il lit des lignes au
clavier, ingère un fichier de données, découpe du texte et le transforme en
nombres. Ce chapitre approfondit les \textbf{flux} (\emph{streams}) de C++ —
\code{std::cin}, \code{std::cout}, et surtout leur cousin discret mais essentiel
\code{std::istringstream} — puis le \textbf{parsing} robuste de texte vers
nombres avec \code{std::stoll} et \code{std::stod}. Si tu viens de Python, oublie
\code{input()} et \code{int(s)} : ici, lire échoue \emph{silencieusement} si tu
ne testes pas l'état du flux, et convertir une chaîne demande de vérifier qu'on
a bien tout consommé. À la fin tu sauras lire jusqu'à la fin d'entrée, découper
une ligne en jetons (\emph{tokens}), et convertir un jeton en nombre sans te
faire avoir par une saisie malformée. Ce chapitre prépare directement le projet
du chapitre~4.
\end{abstract}

\tableofcontents
\bigskip\hrule\bigskip

% ===========================================================================
\section{I/O en profondeur : lire vraiment l'entrée standard}

Tu connais déjà \code{std::cin >> x} et \code{std::cout << x}. On va maintenant
regarder de près \emph{comment} l'entrée se lit, parce que les détails comptent :
sans eux, un programme qui « lit des nombres » se met à boucler à l'infini ou à
ignorer la moitié de la saisie.

\subsection*{Lire mot par mot : \code{std::cin >> mot}}

L'opérateur d'extraction \code{>>} lit \textbf{un jeton} (\emph{token}) à la fois.
Un jeton, c'est une suite de caractères \emph{non blancs}. Point crucial :
\code{>>} \textbf{saute d'abord tous les espaces blancs} (espaces, tabulations,
sauts de ligne) avant de commencer à lire, puis s'arrête au premier blanc
rencontré.

\begin{lstlisting}
#include <iostream>
#include <string>

std::string mot;
std::cin >> mot;   // saute les blancs en tete, lit jusqu'au prochain blanc
\end{lstlisting}

Si l'utilisateur tape \code{\ \ \ bonjour\ tout\ le\ monde}, alors
\code{std::cin >> mot} place \code{"bonjour"} dans \code{mot} (les espaces de tête
sont ignorés, le \code{>>} s'arrête à l'espace après \code{bonjour}). Les mots
suivants restent dans le tampon, prêts pour la prochaine extraction.

C'est exactement pareil pour les types numériques : \code{std::cin >> n} avec
\code{n} un \code{int} saute les blancs puis lit autant de chiffres que possible.

\subsection*{Lire une ligne entière : \code{std::getline}}

Quand tu veux \emph{toute la ligne}, espaces compris, \code{>>} ne convient pas
(il s'arrête au premier espace). On utilise alors la fonction libre
\code{std::getline} :

\begin{lstlisting}
// Signature (simplifiee) :
//   std::istream& std::getline(std::istream& in, std::string& ligne);
// Lit dans 'ligne' tout jusqu'au prochain '\n' (NON inclus dans 'ligne'),
// et consomme ce '\n' dans le flux.

#include <iostream>
#include <string>

std::string ligne;
std::getline(std::cin, ligne);   // lit une ligne complete, espaces compris
\end{lstlisting}

\code{std::getline} lit tout jusqu'au saut de ligne, qu'elle retire du flux mais
\emph{n'ajoute pas} à \code{ligne}. Une ligne vide donne donc une chaîne vide.

\begin{piege}[Mélanger \code{>>} et \code{getline}]
\code{std::cin >> n} lit le nombre mais \emph{laisse} le \code{'\textbackslash n'}
dans le tampon. Si tu enchaînes avec \code{std::getline(std::cin, ligne)}, ce
\code{getline} récupère aussitôt ce saut de ligne résiduel et te rend une ligne
\emph{vide}, sans rien attendre de l'utilisateur. C'est un grand classique. La
parade : ne pas mélanger les deux styles, ou nettoyer le tampon entre les deux
(on verra \code{ignore} au besoin). Pour ce chapitre, choisis un style et
tiens-t'y.
\end{piege}

\subsection*{Lire jusqu'à la fin : la boucle \code{while}}

Comment savoir quand l'entrée est finie ? Au clavier, l'utilisateur signale la
fin d'entrée (\emph{EOF}, \emph{end of file}) avec \textbf{Ctrl-D} sous
Unix/Linux (Ctrl-Z sous Windows). Quand on lit depuis un fichier ou un \emph{pipe}
(\code{./prog < data.txt}), EOF arrive naturellement à la fin du fichier.

L'astuce clé : \textbf{une opération de lecture renvoie le flux lui-même, et un
flux se teste comme un booléen}. Une lecture qui réussit donne « vrai », une
lecture qui échoue (EOF atteint, ou donnée malformée) donne « faux ». D'où les
deux idiomes canoniques :

\begin{lstlisting}
// 1) Lire des valeurs typees jusqu'a EOF ou erreur :
int x;
long long somme = 0;
while (std::cin >> x) {        // s'arrete a EOF ou si la saisie n'est pas un entier
    somme += x;
}

// 2) Lire des LIGNES jusqu'a EOF :
std::string ligne;
while (std::getline(std::cin, ligne)) {   // s'arrete a EOF
    // traiter 'ligne'
}
\end{lstlisting}

La condition \code{while (std::cin >> x)} se lit : « tant que j'arrive à extraire
un entier ». Dès que \code{>>} échoue — parce qu'on est à la fin, ou parce que la
saisie n'est pas un entier — la boucle s'arrête.

\subsection*{Tester et comprendre l'état du flux}

Un flux a quelques \emph{bits d'état} qui expliquent \emph{pourquoi} une lecture
a échoué :

\begin{itemize}[noitemsep]
    \item \code{eof()} : on a atteint la fin de l'entrée.
    \item \code{fail()} : la dernière opération a échoué (ex. on attendait un
          nombre, on a lu une lettre).
    \item \code{good()} : tout va bien, on peut continuer.
\end{itemize}

\begin{lstlisting}
int n;
if (std::cin >> n) {
    std::cout << "lu : " << n << '\n';
} else if (std::cin.eof()) {
    std::cout << "fin d'entree\n";
} else {
    std::cout << "saisie invalide\n";
    // ATTENTION : apres un echec, le flux reste en etat 'fail'.
    // Pour reprendre la lecture il faut std::cin.clear(); puis vider le mauvais token.
}
\end{lstlisting}

\begin{piege}[Un flux en échec reste en échec]
Une fois \code{>>} en échec (\code{fail()} vrai), \emph{toutes} les lectures
suivantes échouent aussi tant que tu n'as pas appelé \code{std::cin.clear()} pour
remettre l'état à zéro. C'est pour ça qu'une boucle \code{while (std::cin >> x)}
qui rencontre une lettre s'arrête net : c'est souvent le comportement voulu, mais
sache-le.
\end{piege}

\begin{aretenir}[Les deux boucles de lecture à connaître par cœur]
\code{while (std::cin >> x)} → lire des \textbf{jetons typés} (entiers, doubles,
mots) un par un, en sautant les blancs, jusqu'à EOF ou première saisie invalide.\\
\code{while (std::getline(std::cin, ligne))} → lire \textbf{ligne par ligne},
espaces compris, jusqu'à EOF.\\
Dans les deux cas, la fin se déclenche au clavier par \textbf{Ctrl-D}, ou
naturellement en fin de fichier avec \code{./prog < fichier}.
\end{aretenir}

% ===========================================================================
\section{\code{std::istringstream} : brancher un flux sur une chaîne}

Voici l'outil central du chapitre. Souvent tu as déjà une \code{std::string} en
main (une ligne lue par \code{getline}, par exemple) et tu veux la \emph{découper}
en morceaux. Plutôt que de manipuler des indices à la main, on \textbf{branche un
flux} sur la chaîne et on réutilise tout ce qu'on vient d'apprendre sur \code{>>}.

C'est le rôle de \code{std::istringstream} (de l'en-tête \code{<sstream>}) : un
flux d'entrée qui lit \emph{depuis une chaîne} au lieu du clavier.

\begin{lstlisting}
#include <sstream>
#include <string>
#include <iostream>

std::string ligne = "12   7  -3   42";
std::istringstream flux(ligne);   // branche un flux sur le contenu de 'ligne'

std::string mot;
while (flux >> mot) {             // meme idiome que std::cin >> mot !
    std::cout << "[" << mot << "]\n";
}
// Affiche [12] [7] [-3] [42] : les espaces multiples sont geres tout seuls.
\end{lstlisting}

Tout ce que tu sais sur \code{std::cin} s'applique : \code{>>} saute les blancs,
s'arrête au prochain blanc, et \code{while (flux >> mot)} s'arrête en fin de
chaîne. \textbf{Les espaces multiples entre les mots sont absorbés
automatiquement} — pas besoin de gérer ce cas toi-même.

\subsection*{Lecture typée : laisser le flux convertir}

Mieux encore : \code{>>} convertit selon le \emph{type} de la variable de
destination. Tu peux donc lire directement des nombres, ou alterner les types.
Imaginons une ligne du genre \code{"3.5 + 2.0"} qu'on veut analyser comme « un
double, un opérateur, un double » :

\begin{lstlisting}
#include <sstream>

std::istringstream flux("3.5 + 2.0");
double a, b;
char op;
if (flux >> a >> op >> b) {        // lit un double, un char, un double
    std::cout << a << ' ' << op << ' ' << b << '\n';   // 3.5 + 2
}
\end{lstlisting}

Chaque \code{>>} saute les blancs puis lit selon le type attendu : \code{a} et
\code{b} sont des \code{double} (le flux convertit le texte en flottant),
\code{op} est un \code{char} (il lit un seul caractère non blanc, ici \code{'+'}).
Le \code{if} global vérifie que les \emph{trois} extractions ont réussi.

\begin{aretenir}[Le réflexe « découper une ligne »]
Tu as une \code{std::string} à analyser ? Branche un \code{std::istringstream}
dessus et lis-la avec \code{>>}, exactement comme \code{std::cin}. Tu profites
gratuitement du saut des blancs, de la gestion des espaces multiples, de la
conversion typée et du test d'état. C'est l'idiome standard pour
\emph{tokenizer} en C++ — bien plus sûr que de découper à la main.
\end{aretenir}

\begin{piege}[\code{>>} sur \code{char} ne lit pas les blancs]
\code{flux >> op} avec \code{op} un \code{char} saute les blancs comme les autres
extractions : il te donne donc le premier caractère \emph{non blanc}. Si tu
voulais lire un espace ou tester la présence exacte d'un séparateur, \code{>>}
n'est pas le bon outil — il faudrait \code{flux.get(c)} qui lit un caractère brut,
blanc compris. À garder en tête, mais pour découper en jetons, \code{>>} est ce
qu'on veut.
\end{piege}

% ===========================================================================
\section{Convertir du texte en nombre, proprement}

Découper une ligne te donne des \code{std::string} comme \code{"42"} ou
\code{"3.14"}. Reste à les transformer en vrais nombres. C++ offre une famille de
fonctions dans \code{<string>} : \code{std::stoi}, \code{std::stoll},
\code{std::stod}\ldots\ (\emph{string to int / long long / double}). On va se
concentrer sur \code{std::stoll} et \code{std::stod}, et surtout sur la manière de
les utiliser \textbf{de façon robuste}.

\subsection*{Les signatures, et l'argument \code{\&pos}}

\begin{lstlisting}
// Signatures (simplifiees) :
//   long long std::stoll(const std::string& s, std::size_t* pos = nullptr, int base = 10);
//   double    std::stod (const std::string& s, std::size_t* pos = nullptr);

#include <string>

std::string s = "42";
long long n = std::stoll(s);   // n vaut 42
double    x = std::stod("3.14");  // x vaut 3.14
\end{lstlisting}

Appelées simplement, ces fonctions convertissent le \emph{début} de la chaîne et
\textbf{ignorent ce qui suit}. C'est le piège majeur : \code{std::stoll("42xyz")}
renvoie \code{42} sans broncher, et \code{std::stoll("12 34")} renvoie \code{12}.
Pour beaucoup d'usages, « tolérer du rab » est inacceptable.

D'où le deuxième paramètre, \code{pos} : un \emph{pointeur} vers un
\code{std::size\_t} que la fonction \textbf{remplit avec l'indice du premier
caractère non consommé}. Après l'appel, \code{pos} vaut le nombre de caractères
lus. On le compare à la taille de la chaîne pour vérifier qu'on a tout consommé :

\begin{lstlisting}
#include <string>

std::string s = "42xyz";
std::size_t pos = 0;
long long n = std::stoll(s, &pos);   // n == 42, pos == 2 (s'est arrete a 'x')

if (pos == s.size()) {
    // toute la chaine a ete consommee : conversion stricte reussie
} else {
    // il reste des caracteres non numeriques apres le nombre
}
\end{lstlisting}

On passe \code{\&pos} (l'\emph{adresse} de \code{pos}) parce que le paramètre est
un pointeur \code{std::size\_t*} : la fonction écrit \emph{à travers} ce pointeur
pour nous renvoyer où elle s'est arrêtée.

\subsection*{Les exceptions : \code{invalid\_argument} et \code{out\_of\_range}}

Contrairement à \code{int(s)} en Python qui lève toujours, ces fonctions
\textbf{lèvent une exception} dans deux cas, qu'il faut attraper :

\begin{itemize}[noitemsep]
    \item \code{std::invalid\_argument} : \textbf{aucune} conversion possible —
          la chaîne ne commence pas par un nombre (ex. \code{std::stoll("abc")}
          ou \code{std::stoll("")}).
    \item \code{std::out\_of\_range} : un nombre a bien été lu, mais il
          \textbf{dépasse} la capacité du type cible (ex. un entier gigantesque
          pour \code{stoll}).
\end{itemize}

\begin{lstlisting}
#include <string>
#include <stdexcept>

try {
    std::size_t pos = 0;
    long long n = std::stoll("abc", &pos);   // lance std::invalid_argument
    (void) n;
} catch (const std::invalid_argument&) {
    // rien de numerique au debut
} catch (const std::out_of_range&) {
    // trop grand pour un long long
}
\end{lstlisting}

\subsection*{Un parse strict, tout assemblé}

Voici le motif complet qu'on réutilisera partout. Il combine les trois
vérifications : exceptions \emph{et} consommation totale via \code{pos}.

\begin{lstlisting}
#include <string>
#include <stdexcept>
#include <optional>

// Renvoie le nombre si la chaine ENTIERE est un long long valide, sinon "rien".
std::optional<long long> parse_strict(const std::string& s) {
    try {
        std::size_t pos = 0;
        long long n = std::stoll(s, &pos);
        if (pos != s.size()) {
            return std::nullopt;   // des caracteres en trop apres le nombre
        }
        return n;                  // OK : toute la chaine consommee
    } catch (const std::invalid_argument&) {
        return std::nullopt;       // rien de numerique
    } catch (const std::out_of_range&) {
        return std::nullopt;       // hors capacite
    }
}
\end{lstlisting}

Note un détail subtil : \code{std::stoll(" 42")} (avec un espace de tête) réussit
et renvoie \code{42}, car la fonction saute les blancs initiaux — mais un blanc
\emph{en fin} (\code{"42 "}) fait que \code{pos != s.size()}. Selon ton besoin, tu
choisiras d'ôter les blancs autour avant, ou non. L'important est de
\emph{décider} consciemment ce qui est valide.

\begin{aretenir}[\code{istringstream} pour découper vs \code{stoll}/\code{stod} pour convertir]
Ce sont deux outils complémentaires, ne les confonds pas :
\begin{itemize}[noitemsep]
    \item \code{std::istringstream} \textbf{découpe} : on lui donne une ligne
          \emph{entière} contenant plusieurs jetons, et \code{>>} la débite mot
          par mot en sautant les blancs. C'est l'outil du \emph{tokenizing}.
    \item \code{std::stoll} / \code{std::stod} \textbf{convertissent} : on leur
          donne \emph{un seul jeton déjà isolé} (\code{"42"}) et ils en font un
          nombre, avec contrôle strict via \code{\&pos} et gestion des
          exceptions.
\end{itemize}
Le pipeline typique : \code{getline} lit une ligne → \code{istringstream} la
découpe en jetons → \code{stoll}/\code{stod} convertit chaque jeton. À noter :
\code{flux >> n} (lecture typée directe) fait découpage \emph{et} conversion d'un
coup, mais signale l'échec par l'état du flux plutôt que par une exception — pas
par \code{pos}. Garde \code{stoll}/\code{stod} quand tu veux un contrôle fin sur
un token isolé venu d'ailleurs.
\end{aretenir}

\begin{piege}[\code{std::stoll} accepte un signe et une base]
\code{std::stoll("-7")} vaut \code{-7} (le signe est permis), et le 3\textsuperscript{e}
argument \code{base} permet de lire de l'hexadécimal (\code{std::stoll("ff", \&pos, 16)}
vaut 255). Pratique, mais ça veut dire que \code{"-7"} ou \code{"+3"} passent ton
parse « entier » — vérifie que c'est bien ce que tu veux pour ton exercice.
\end{piege}

% ===========================================================================
\section*{À toi de jouer}
\addcontentsline{toc}{section}{À toi de jouer}

\bigskip\hrule\bigskip
\noindent Les exercices du dossier \code{exercices/} font travailler chaque
brique de ce chapitre, en difficulté croissante. Tu y écriras toi-même la
logique ; l'énoncé rappelle au besoin l'outil utile.

\begin{itemize}[noitemsep]
    \item \textbf{Parsing de lignes} : lire l'entrée ligne par ligne avec
          \code{getline}, brancher un \code{istringstream} sur chaque ligne, et
          en extraire les champs voulus. Le pain quotidien du traitement de
          données.
    \item \textbf{Tokenize isolé} : écrire une fonction qui prend une
          \code{std::string} et renvoie un \code{std::vector<std::string>} de ses
          jetons, en t'appuyant sur \code{istringstream}. Brique réutilisable
          partout ensuite.
    \item \textbf{Conversion stricte} : implémenter un \code{parse} robuste
          (\code{stoll}/\code{stod} + \code{\&pos} + exceptions) qui renvoie un
          \code{std::optional} — « le nombre, ou rien si la chaîne est invalide ».
    \item \textbf{Lecture-stats} : programme \emph{autonome} (\code{main}
          compris) qui lit une suite de nombres sur l'entrée standard jusqu'à
          EOF et affiche compte, somme, min, max, moyenne. L'occasion de maîtriser
          \code{while (std::cin >> x)} et le test d'état.
    \item \textbf{Mini-calculatrice} : programme \emph{autonome} qui lit des
          lignes du type \code{a op b} (\code{double char double}), évalue, et
          gère proprement les saisies invalides. Tu y combines lecture typée,
          \code{istringstream} et gestion d'erreurs.
\end{itemize}

\begin{aretenir}[Cap sur le chapitre 4]
Ces réflexes — lire jusqu'à EOF, découper avec \code{istringstream}, convertir
strictement avec \code{stoll}/\code{stod} — sont exactement ceux dont tu auras
besoin pour le \textbf{projet du chapitre~4}, qui assemblera tout ça dans une
application qui lit, parse et traite un vrai flux de données. Soigne ces bases :
elles portent tout le reste.
\end{aretenir}

\end{document}
